\section{Discussion}
\label{sec:discussion}

The paper's findings are two sides of one mechanism. The lifted system's dominant error is
\emph{over-counting}: one physical object survives cross-view fusion as several same-label nodes
(Table~\ref{tab:perclass_overall}). As views are added, recall rises but precision falls faster, so
accuracy is best at the sparse end. Over-counting is an over-\emph{splitting} error, so the
corrective move is to \emph{merge} --- and the move in the opposite direction, an uncertainty gate
that \emph{tightens} merging, can only make it worse. This explains both empirical results at once.

\subsection{Why merging helps and the uncertainty gate does not}
\label{sec:disc-mechanism}

Duplicate-aware fusion merges same-label nodes whose axis-aligned 3D boxes overlap, directly
targeting the over-counting. It recovers precision on exactly the over-counted classes while
preserving recall (Table~\ref{tab:dedup_mechanism}), and is the only variant whose accuracy
\emph{rises} with more views --- each added view supplies redundancy the merge absorbs rather than
fragmentation it cannot. The uncertainty-aware fusion gate acts on the same axis in the wrong
direction: it tightens merging precisely where the system already over-splits. No weight in the
sweep $w \in [0.1, 0.8]$ beats the non-gated baseline, and the default regresses by $-0.033$ to
$-0.096$ across the four view counts (Table~\ref{tab:uncertainty_neg}). This is not a tuning failure
but the predictable cost of suppressing merges in a regime whose error is too few merges; the
contrast between the two variants is itself the clearest evidence for the over-counting diagnosis.

\subsection{The circular-evaluation lesson}
\label{sec:disc-circularity}

This mechanism was invisible under a prediction-seeded reference. When the reference is seeded from
the model's own dense-view output, a method that reproduces its ten-view graph scores $1.0$ at the
reference view \emph{by construction}, and recall-raising behavior is rewarded whether or not it
corresponds to real objects (Section~\ref{subsec:decirc}); the uncertainty gate then looked like a
sparse-view win. Re-grading the identical outputs against an independently authored reference
reverses that result and exposes the precision collapse underneath. The takeaway generalizes: a
reference seeded from a model measures agreement with that model, not recognition, and silently
awards a perfect score to whatever produced the seed --- an easy trap in sparse-view 3D, where dense
ground truth is costly~\cite{gundersen2018reproducibility,kapoor2023leakage}. We therefore treat the
independent reference and the documented pitfall as a first-class contribution, and recommend
reporting any prediction-seeded reference alongside an independent one.

\subsection{Limitations}
\label{sec:disc-limitations}

\paragraph{Reference provenance.} Of the 28 scenes, 5 are human-verified (the gold core) and 23
carry two-pass VLM-drafted references pending human verification. Both findings reach $p<0.001$ over
the 28 scenes and hold on the gold five alone (Section~\ref{subsec:dedup}), so they do not rest on
the drafted labels; fully verifying the 23 is the natural next step, and the benchmark is small
relative to large posed-RGB-D datasets.

\paragraph{Closely-packed same-class instances.} Box-IoU merging cannot separate distinct same-class
objects whose axis-aligned 3D boxes overlap, such as two monitors on one desk. This is the merge's
only recall cost --- $-0.04$ on \texttt{desk} and \texttt{bottle} at eight views, with other classes
unchanged (Table~\ref{tab:dedup_mechanism}) --- and the precision gain dwarfs it.

\paragraph{Static, ``thing''-centric, rule-based.} The pipeline assumes a static scene, so VGGT's
multi-view geometry degrades on the dynamic (moving-person) sequences. ``Stuff'' classes
(\texttt{wall}, \texttt{floor}, \texttt{ceiling}) are treated as countable objects, a poor fit for
amorphous surfaces. And the spatial relations come from centroid/distance rules that are auditable
but cannot express the support, containment, or functional relations of richer 3D scene
graphs~\cite{armeni2019scenegraph,wald2020threedssg,gu2024conceptgraphs}.

\subsection{Future work}
\label{sec:disc-future}

The merge's one failure mode points to the fix: appearance-based instance separation that splits a
merged group when its members are visually distinct. This is also where the recall signal the
uncertainty gate tried to add belongs --- not as a merge gate that blocks fusion in the
over-counting regime, but as a \emph{split} decision on an already-merged group, the one operation
where a recall signal helps rather than hurts. Replacing the rule-based predicates with a learned
relation module, once the independent annotations supply enough supervision, is the other obvious
extension.

\section{Conclusion}
\label{sec:conclusion}

We presented a detector-grounded open-vocabulary 3D scene-graph system that turns a handful of casual
RGB views into an inspectable object--relation graph with all heavy models frozen and no per-scene
optimization: OWLv2~\cite{minderer2023owlv2} proposes labeled boxes, feed-forward
VGGT~\cite{wang2025vggt} geometry lifts them to 3D, and cross-view fusion assembles the graph. On an
independent TUM RGB-D~\cite{sturm2012tumrgbd} benchmark, the dominant error is over-counting, and a
one-line duplicate-aware merge corrects it, improving object-label F1 by $+0.06$ to $+0.12$ at every
view count and winning on $22$--$27$ of the $28$ scenes ($p<0.001$); a recalibrated, fully swept
uncertainty-aware gate --- the opposite move --- gives no benefit, a negative result that a
prediction-seeded reference had disguised as a win. The de-circularized reference and the documented
pitfall are the methodological core; the benchmark and protocol are dataset-agnostic, and we release
them for reuse.
